%%
%% Supplementary material for the Bioinformatics Applications Note.
%% Tables under tables/ are generated by make_supplementary_tables.py from
%% benchmarks/results/paper_data.json, so they cannot drift from the artifacts.
%%
\documentclass[11pt]{article}

\usepackage[T1]{fontenc}
\usepackage[utf8]{inputenc}
\usepackage{lmodern}
\usepackage[margin=1in]{geometry}
\usepackage{booktabs}
\usepackage{graphicx}
\usepackage{amsmath}
\usepackage{xcolor}
\usepackage{listings}
\usepackage[hidelinks]{hyperref}
\usepackage{caption}

\captionsetup{font=small,labelfont=bf}
\lstset{
  basicstyle=\ttfamily\footnotesize,
  breaklines=true,
  frame=single,
  columns=fullflexible,
  backgroundcolor=\color{black!3},
}
\newcommand{\code}[1]{\texttt{#1}}

\title{\bfseries Supplementary material\\[2pt]
\large lit-acquisition: multi-source, multilingual biomedical literature
acquisition with citation-graph traversal}
\author{}
\date{}

\begin{document}
\maketitle
\vspace{-3em}

\section{Reference set: construction and what it can support}
The reference set contains 43 clinically actionable variants and 822 papers.
Construction followed a documented automated protocol: for each variant a
clinical-style natural-language query was formulated, candidates were pooled
from PubMed E-utilities, Europe PMC and Crossref, and a weighted heuristic
scored each candidate (gene symbol in title $+3$; variant identifier $+5$;
disease keywords up to $+3$; case-report publication type $+2$; reviews $-2$;
editorials, letters, comments and news excluded). The 30 original variants were
annotated in two rounds; 13 further variants were added on 2026-08-14 by
\code{benchmarks/gold\_standard/build\_variants.py}, which replicates the same
protocol without modifying the originals.

\textbf{Limitations that bound every claim in the paper.}
(i) The annotations are \emph{heuristic and not expert-verified}, so they
contain unquantified false positives and false negatives.
(ii) The candidate pool was drawn from channels that \emph{overlap the evaluated
system}, a pooling/circularity bias: absolute recall against this set measures
agreement with the annotation protocol, not clinical completeness.
(iii) Language labels come from provider metadata; 305 of 822 papers have no
label, and only German (13) and French (11) carry non-English annotations, so
claims about other target languages are not supported by this set.
Consequently all reported metrics are \emph{relative} coverage comparisons.
\code{benchmarks/gold\_standard/variants.jsonl} is the source of truth;
\code{variants\_annotated.jsonl} is a historical V001--V030 file.

\section{Configurations and evaluation protocol}
Four configurations were compared at a uniform budget of 100 keyword candidates
per variant:
\begin{itemize}\itemsep2pt
  \item \textbf{M1 PubMed-only} --- PubMed E-utilities only, English query.
  \item \textbf{M2 multi-source} --- all channels, English query.
  \item \textbf{M3 +translation} --- adds LLM translation into en/zh/ja/de/fr/ru
        with language-aware routing.
  \item \textbf{M4n +citation (nested)} --- strictly nested ablation built from
        the frozen M3 ranked list plus citation-graph items recorded in the
        independent M4 run, deduplicated and appended in traversal order
        (monotonicity verified; the keyword phase is not re-executed).
\end{itemize}
Metrics are macro means over the 43 variants: SetRecall (fraction of a variant's
reference papers retrieved), Recall@$K$, Precision@$K$ (fixed denominator $K$),
Precision@Retrieved, and nDCG@$K$ for $K\in\{5,10,15,20,50\}$. Micro aggregates
are also reported in the artifacts. Statistics: two-sided paired Wilcoxon
signed-rank tests with zeros dropped, Holm--Bonferroni correction over a
pre-declared family of eight tests, matched-pairs rank-biserial effect size, and
10{,}000 variant-level percentile bootstrap resamples (seed 42). Live-API
run-to-run variance was estimated from two independent executions of the M3
configuration (18/30 identical ranked lists; mean Jaccard 0.891).

\textbf{Reproducibility of the LLM components.} The recorded collection run used
\code{Qwen/Qwen2.5-72B-Instruct} for translation (temperature 0.2,
max\_tokens 1024, system-prompt SHA-256 prefix \code{96fe414ab84606dd}) and
\code{BAAI/bge-reranker-v2-m3} as the cross-encoder rerank stage. All of these
are recorded in \code{benchmarks/results/run\_manifest.json}. The relevance gate
was \emph{disabled} during collection, so the reported metrics are raw retrieval
comparisons; the lite edition performs no LLM calls at all.

\section{Full metrics}
\begin{table}[h]
\centering\small
\caption{Full per-method metrics over 43 variants (macro means). Precision@$K$ uses a fixed denominator $K$; Precision@Retr.\ is over the retrieved list.}
\label{tab:supp-metrics}
\begin{tabular}{lccccccc}
\toprule
Method & SetRec. & Rec@5 & Rec@20 & Rec@50 & P@20 & nDCG@20 & nDCG@50 \\
\midrule
M1 PubMed-only & 0.269 & 0.076 & 0.171 & 0.218 & 0.155 & 0.205 & 0.219 \\
M2 multi-source & 0.472 & 0.044 & 0.153 & 0.275 & 0.147 & 0.164 & 0.222 \\
M3 +translation & 0.465 & 0.085 & 0.276 & 0.446 & 0.265 & 0.301 & 0.381 \\
M4n +citation (nested) & 0.467 & 0.085 & 0.276 & 0.446 & 0.265 & 0.301 & 0.381 \\
\bottomrule
\end{tabular}
\end{table}


\section{Significance tests}
\begin{table}[h]
\centering\small
\caption{Pre-declared eight-test family: paired Wilcoxon signed-rank with Holm--Bonferroni correction, 10{,}000 bootstrap resamples (seed 42). $r$ is the matched-pairs rank-biserial correlation.}
\label{tab:supp-tests}
\begin{tabular}{lcccc}
\toprule
Contrast & $\Delta$ & 95\% CI & $p_{\mathrm{Holm}}$ & $r$ \\
\midrule
M2 vs M1 set recall & +0.203 & [+0.128, +0.277] & 1.08e-04 & +0.77 \\
M3 vs M2 set recall & -0.007 & [-0.087, +0.069] & 1.00e+00 & +0.02 \\
M4n vs M3 set recall & +0.002 & [+0.000, +0.004] & 1.00e+00 & +1.00 \\
M4n vs M1 set recall & +0.198 & [+0.123, +0.269] & 1.93e-04 & +0.74 \\
M4n vs M2 set recall & -0.005 & [-0.086, +0.071] & 1.00e+00 & +0.03 \\
M4n vs M1 recall20 & +0.105 & [+0.035, +0.170] & 1.61e-02 & +0.51 \\
M4n vs M1 ndcg50 & +0.162 & [+0.086, +0.234] & 4.30e-04 & +0.66 \\
M4n vs M3 prec retr & -0.011 & [-0.016, -0.006] & 1.91e-06 & -1.00 \\
\bottomrule
\end{tabular}
\end{table}


\section{Comparison with single-channel baselines}
\begin{table}[h]
\centering\small
\caption{Single-channel baselines (uniform 20-paper budget over the same PubMed candidate pool unless noted). Significance is versus M4n.}
\label{tab:supp-sota}
\begin{tabular}{lccc}
\toprule
Method & SetRecall & Recall@20 & $p_{\mathrm{Holm}}$ vs M4n \\
\midrule
S1 Dense (bge-m3) & 0.269 & 0.182 & 3.86e-04 \\
S2 Cross-encoder rerank & 0.269 & 0.171 & 3.86e-04 \\
S3 Hybrid BM25+dense+rerank & 0.421 & 0.271 & 1.00e+00 \\
S4 LLM query expansion & 0.270 & 0.173 & 3.86e-04 \\
S5 Agentic 2-round search & 0.269 & 0.171 & 3.86e-04 \\
S6 Multilingual dense (bge-m3) & 0.269 & 0.181 & 3.86e-04 \\
\midrule
\textbf{M4n pipeline} & \textbf{0.467} & \textbf{0.276} & --- \\
\bottomrule
\end{tabular}
\end{table}


Note that the four baselines that only reorder the shared PubMed candidate pool
(mean 80.7 candidates: dense, cross-encoder, agentic and multilingual dense)
\emph{cannot} exceed M1 on set recall by construction --- their SetRecall
columns are identical to M1's. The informative comparison is the hybrid
BM25+dense+cross-encoder baseline.

\section{Multilingual recovery}
\begin{table}[h]
\centering\small
\caption{Annotated German/French reference papers recovered per method. Only these two languages carry non-English annotations.}
\label{tab:supp-language}
\begin{tabular}{lccc}
\toprule
Method & German (of 13) & French (of 11) & Total (of 24) \\
\midrule
M1 PubMed-only & 0 & 0 & 0 \\
M2 multi-source & 1 & 0 & 1 \\
M3 +translation & 11 & 10 & 21 \\
M4n +citation (nested) & 11 & 10 & 21 \\
\midrule
Best English-only baseline (S3 hybrid) & 1 & 0 & 1 \\
\bottomrule
\end{tabular}
\end{table}


Translation is set-level neutral in aggregate (M3 vs M2:
$p_{\mathrm{Holm}}=1.0$) but changes which papers are reachable, which is the
practical argument for language-aware routing.

\section{Citation-graph exclusivity}
Across all 43 variants, citation traversal contributed exactly two relevant
reference papers, one of which appears in no keyword list:
\begin{itemize}\itemsep2pt
  \item V002 (BRCA1 c.5266dupC) --- \emph{Familial history and prevalence of
        BRCA1, BRCA2 and TP53 pathogenic variants in HBOC Brazilian patients}
        (Sci Rep 2022, \code{10.1038/s41598-022-23012-3}).
  \item V035 (HFE p.C282Y) --- \emph{HFE p.C282Y Polymorphism and Risk of
        Metabolic Syndrome Components: Systematic Review and Meta-Analysis}
        (Medicina 2026, \code{10.3390/medicina62030589}).
\end{itemize}
The same step costs list precision (Precision@Retrieved $-0.011$,
$p_{\mathrm{Holm}}=1.9\times10^{-6}$), so we present traversal as a bounded
complementary channel rather than a headline gain.

\section{Provider failures and operational notes}
\begin{table}[h]
\centering\small
\caption{Provider failures recorded during the collection session (window 2026-08-12 12:09:28 to 2026-08-12 13:47:18).}
\label{tab:supp-failures}
\begin{tabular}{lc}
\toprule
Signal & Count \\
\midrule
Semantic Scholar HTTP 429 & 78 \\
HTTP 429 (all providers) & 78 \\
J-STAGE parse failures & 120 \\
Citation-graph failures & 0 \\
Translation calls logged & 60 \\
\bottomrule
\end{tabular}
\end{table}


The J-STAGE parse failures were traced to a third-party client that emitted
percent-encoded URLs; the channel was replaced by a pure-Python backend and
returns real results from v0.4.1 onward. Semantic Scholar anonymous access is
aggressively rate-limited from shared/data-centre IPs, which is why an API key
is recommended for batch use.

\section{Cost measurements}
\noindent
\begin{tabular}{@{}lll@{}}
\toprule
Measurement & Value & Artifact \\
\midrule
Full pipeline (composed), median & 28.65\,s/variant & \code{timing\_43variants.json} \\
Full pipeline (composed), p90 & 44.61\,s/variant & \code{timing\_43variants.json} \\
Download rate (end-to-end) & 3.99 PDFs/min & \code{download\_throughput.json} \\
Candidate$\rightarrow$PDF conversion & 9.6\% (18/187) & \code{download\_throughput.json} \\
Seconds per PDF & 15.06 & \code{download\_throughput.json} \\
Peak RSS & 85.7\,MB & \code{download\_throughput.json} \\
\bottomrule
\end{tabular}

\medskip
The timing figure is a composition of the translation-phase run and the
citation phase from the same 43-variant session, not a fresh end-to-end run, so
that it stays tied to the frozen benchmark lists. The download measurement
covers four variants and is a small-sample estimate; it also showed that the
download layer returns nothing without \code{LIT\_UNPAYWALL\_EMAIL}, because
every DOI$\rightarrow$open-access route is then skipped.

\section{Reproducibility}
\begin{lstlisting}
pip install lit-acquisition              # lite edition (no LLM calls)
pip install "lit-acquisition[agent]"     # + translation, gate, rerank

python benchmarks/run_benchmark.py --collect --limit 100 --cg-budget 20
python benchmarks/make_paper_data.py            # regenerates paper numbers
python benchmarks/measure_download_and_memory.py --variants 4 --limit 8
python benchmarks/compose_timing.py
python paper/appnote/make_figure1.py
python paper/appnote/make_supplementary_tables.py
python paper/appnote/check_numbers.py           # verifies the manuscript
\end{lstlisting}
Environment of the recorded run: Python 3.12.3, Linux. Provider API keys were
read from the environment; the Unpaywall contact email is required for the
DOI$\rightarrow$open-access route. Reference-set and result artifacts live under
\code{benchmarks/}; they are excluded from the PyPI distribution and must be
taken from the repository.

\section{Usage examples}
\noindent\textbf{Search across all channels (lite edition).}
\begin{lstlisting}
import asyncio
from lit_acquisition import online_acquisition_workflow

async def main():
    result = await online_acquisition_workflow({
        "query": "MECP2 Rett syndrome case report",
        "action": "search",
        "limit": 20,
    })
    print(result["summary"])          # one-line outcome
    print(result["diagnostics"])      # per-provider status/latency/counts

asyncio.run(main())
\end{lstlisting}

\noindent\textbf{Multilingual acquisition with downloads (agent edition).}
\begin{lstlisting}
from lit_acquisition import multilingual_acquisition_workflow

result = await multilingual_acquisition_workflow({
    "query": "HFE p.C282Y hereditary haemochromatosis",
    "action": "download",
    "language": "multilingual",
    "limit": 30,
    "download_path": "./downloads",
})
\end{lstlisting}

\noindent\textbf{Citation-graph traversal from a seed DOI.}
\begin{lstlisting}
from lit_acquisition import traverse_citation_graph

papers = await traverse_citation_graph(
    seed="10.1038/s41598-022-23012-3", max_depth=1, max_papers=20, direction="both",
)
\end{lstlisting}

\end{document}
